\section{The fluctuation covariance and the step}\label{sec:covariance}

We return to the setting of Section~\ref{sec:fluctuations}: $W=V^0$ is the base configuration of the step, $U'=U_{k+1}(V')$, and $\eps=B_3\eps'_1$. The objects of Section~\ref{sec:coarse-action} are taken at this $W$ and $U'$.

\subsection{The quadratic form of the step}

\begin{lemma}[Expansion in the free variables]\label{lem:step-quadratic}
Let $r_2=\min(cL^{-2d-2},\,cL^{-d}a'_1)$ with $c>0$ small, depending only on $d$, $G$, $L$ and $r_1=c_6/4$ of Proposition~\ref{prop:coarse-action} (the proof needs $2cL^{-2d-2}\le\frac12r_L$, $CL^{-d-2}c\le\frac14r_1$ and $Cc\le\frac12$), $a'_1$ as in Lemma~\ref{lem:composition}(c), and assume $4(C_*+1)L^{-2}\eps'_1\le a'_1$. For complex fields $Z$ on $\mathfrak F$ with $|Z|_\infty\le r_2$ put $F(Z)=\mathcal S(\hat Z(Z))$. Then $F$ is analytic, $F(Z)=S(U_k(e^{\ii\hat Z(Z)}W))$ for real $Z$, and
\[
F(Z)=S(U')+\tfrac12\Delta_k(Z,Z)+\mathcal V(Z),\qquad\Delta_k(Z,Z)=\Delta^0(\hat Z^1Z,\hat Z^1Z)+2\ip{\omega_0}{\hat Z^2(Z)},
\]
with $|\mathcal V(Z)|\le C_L|Z|_\infty\|Z\|^2$. In particular $F$ has no linear term.
\end{lemma}

\begin{proof}
By Lemma~\ref{lem:eliminate}, $\hat Z$ is analytic on $|Z|_\infty\le2r_2\le r_L$ with $|\hat Z(Z)|_\infty\le CL^d|Z|_\infty\le r_1/2$, so $F$ is analytic on the disc of radius $2r_2$ by Proposition~\ref{prop:coarse-action}(a). For real $Z$ the configuration $e^{\ii\hat Z(Z)}W$ satisfies the one-step constraint, and its curvature is at most $2(C_*+1)L^{-2}\eps'_1+CL^dr_2\le a'_1$. Here $|\hat Z|_\infty\le CL^d|Z|_\infty$, so $CL^dr_2\le a'_1/2$ by the choice of $r_2$, and we assume $4(C_*+1)L^{-2}\eps'_1\le a'_1$. By Lemma~\ref{lem:composition}(c), $F(Z)\ge F(0)=S(U')$, so $dF(0)=\ip{\omega_0}{\hat Z^1\cdot}=0$. The second-order term follows from $\hat Z(tZ)=t\hat Z^1Z+t^2\hat Z^2(Z)+O(t^3)$ and Proposition~\ref{prop:coarse-action}(d). For the remainder, $\nabla F(Z)=d\hat Z(Z)^T\omega(\hat Z(Z))$. The analyticity of $\hat Z$ and the bound $|\hat Z(Z)|_\infty\le r_1/2$ hold on the larger disc $|Z|_\infty\le2cL^{-2d-2}$, since only the first term of the minimum defining $r_2$ enters them. By Proposition~\ref{prop:coarse-action}(c) and the locality of $\hat Z$ (range $2L$), $d^2F$ has a kernel bounded by $C_Le^{-\gamma_1|b-b'|}$ on this larger disc. The argument of Proposition~\ref{prop:coarse-action}(d), applied to $F$ on this disc with $dF(0)=0$, gives the bound on $\mathcal V$. Its constant $C_L$ depends on $cL^{-2d-2}$, but not on $a'_1$.
\end{proof}

This is the quadratic form \cite[(53)--(54)]{BalUV3}. The term $2\ip{\omega_0}{\hat Z^2}$ comes from the curvature of the constraint set and is Ba\l aban's term $-\ip{H_1\hat D^{(2)}(A)}J$.

\subsection{Coercivity}

The lower bound for $\Delta_k$ is proved in two steps:
\begin{itemize}
\item at the trivial background, a Poincar\'e inequality on the one-step constraint set combined with Proposition~\ref{prop:eff-action-gauge};
\item at a regular background, an IMS localization on the unit lattice, as in Section~\ref{sec:reg-gauge}.
\end{itemize}
Let $q$ be the one-step tube average at the trivial background, i.e.\ $l'$ with $W=1$.

\begin{lemma}[Poincar\'e inequality for one step]\label{lem:poincare-step}
There is $C_{\rm P}$, depending only on $d$ and $L$, such that every real field $B$ on the bonds of $\T$ with $B=0$ on $\mathcal T$ satisfies
\[
\|B\|^2\le C_{\rm P}\big(\|\dl B\|^2+\|qB\|^2\big),
\]
where $\dl B$ is the unit-lattice curl.
\end{lemma}

\begin{proof}
All estimates are local in a pair of neighbouring blocks.

\emph{Bonds inside a block.} Let $b=(x,x')$ be a bond of $B'(y)$ not in $\mathcal T$. By Lemma~\ref{lem:stokes}(ii) the loop $\Gamma^{k+1}_{y,x}\cup b\cup(\Gamma^{k+1}_{y,x'})^{-1}$ is spanned by at most $dL$ plaquettes. Since $B=0$ on the tree, the abelian Stokes formula gives $B(b)=\sum_i\pm\dl B(p_i)$, so $|B(b)|^2\le dL\sum_i|\dl B(p_i)|^2$.

\emph{Crossing bonds.} Let $c=(y,y')$, $y'=y+Le_\mu$, and for $x$ in the face $\mathsf F=\{x\in B'(y):x_\mu=y_\mu+L-1\}$ let $b_x=(x,x+e_\mu)$. For $x,\tilde x\in\mathsf F$ the closed contour
\[
\Gamma^{k+1}_{y,x}\,b_x\,(\Gamma^{k+1}_{y',x+e_\mu})^{-1}\,\Gamma^{k+1}_{y',\tilde x+e_\mu}\,b_{\tilde x}^{-1}\,(\Gamma^{k+1}_{y,\tilde x})^{-1}
\]
lies in a box of side $2L$ and has at most $4dL+2$ bonds. By Lemma~\ref{lem:stokes}(iii) it is spanned by at most $12d^2L^2$ plaquettes, so $B(b_x)-B(b_{\tilde x})$ is a sum of at most $12d^2L^2$ values $\pm\dl B(p)$.

\emph{The average.} $(qB)(c)=L^{-d}\sum_{x\in B'(y)}L^{-1}B([x,x(c)])$. Each straight contour $[x,x(c)]$ contains exactly one crossing bond $b_{\tilde x}$, with $\tilde x\in\mathsf F$ on the same line as $x$, and $L-1$ bonds inside the two blocks. Each $\tilde x$ occurs for $L$ values of $x$. Hence the mean $m_c$ of $B(b_{\tilde x})$ over $\tilde x\in\mathsf F$ satisfies $|m_c|\le L|(qB)(c)|+(L-1)\max|B(b)|$, the maximum running over the inner bonds of the two blocks. Then $|B(b_x)|\le|m_c|+\max_{\tilde x}|B(b_x)-B(b_{\tilde x})|$.

Squaring and summing, each plaquette and each value $(qB)(c)$ is used a number of times bounded in terms of $d$ and $L$.
\end{proof}

\begin{lemma}[Local comparison]\label{lem:local-comparison}
Let $M'\ge\max(M,L)$ be an integer and $z$ a point of the torus (not necessarily of $\T$). Assume $n\ge80M'$ and $\eps\le\eps_{\rm cmp}(M')$, where $\eps_{\rm cmp}(M')>0$ depends only on $M'$, $d$, $L$, $G$, $a$, $\alpha$ and $M$, and is fixed in the proof. There are a real gauge transformation $u$ on $\Lam_{40M'}(z)$ and $C_{M'}$ such that, for every real field $B$ on the bonds of $\T$ starting in $\Lam_{M'}(z)$ and vanishing on $\mathcal T$, with $B^u=R(u)B$:
\[
\big|\ip B{(\boldsymbol\Gamma_\pi-a)B}-\ip{B^u}{\bD B^u}\big|\le\big(C_{M'}\eps+Ce^{-\gamma M'}\big)\|B\|^2,\qquad\big|\|l'B\|-\|qB^u\|\big|\le C_{M'}\eps\|B\| .
\]
Here $\bD=\boldsymbol\Gamma(1)-a$ is the operator of Proposition~\ref{prop:eff-action-gauge}, and $C,\gamma$ do not depend on $M'$.
\end{lemma}

\begin{proof}
The configuration $U'$ is $3\eps$-regular on every cube $\Lam_{4M+2}(z'')$ (Section~\ref{sec:local-gauges}). Let $z_0\in\Tf$ with $|z-z_0|_\infty\le\eta$, so that $\Lam_{40M'}(z)\subset\Lam_{40M'+1}(z_0)$. Apply Lemma~\ref{lem:large-cubes} with $r=4M+2\ge18$, $R=40M'+1$ and $\sigma=3\eps$, on the cube $\Lam_R(z_0)$. Its hypothesis $C_E(R/r)^pR\sigma\le1$ is the first condition defining $\eps_{\rm cmp}(M')$. The boxes in its proof have sides at most $3R=120M'+3<80M'L\le Ln$, the side of $\T$, so they do not wrap around the torus. This gives a real gauge $u$ on $\Lam_{40M'}(z)$ with $U'^u=e^{\ii\eta A}$ and $\|A\|_{1,\alpha}\le C(40M'/M)^p\eps$. Extend $u$ to a real gauge transformation $\tilde u$ of $\Tf$. Then $\boldsymbol\Gamma_\pi(U'^{\tilde u})=R(\tilde u)\boldsymbol\Gamma_\pi(U')R(\tilde u)^{-1}$ and $l'(W^{\tilde u})R(\tilde u)=R(\tilde u)l'(W)$, by Proposition~\ref{prop:linearized-averages}(f), Proposition~\ref{prop:covariance} and Lemma~\ref{lem:one-step-basic}(b).

Let $\chi$ be a smooth cutoff equal to $1$ on $\Lam_{20M'}(z)$ and $0$ outside $\Lam_{30M'}(z)$, and put $U^\natural=e^{\ii\eta\chi A}$. It satisfies (B) with parameter $C_{M'}\eps$. The second condition defining $\eps_{\rm cmp}(M')$ is that $C_{M'}\eps$ lies below the thresholds of Parts~II and~III. The configuration $U^\natural$ coincides with $U'^{\tilde u}$ on $\Lam_{20M'}(z)$. By Theorem~\ref{thm:reg-locality}(b), the kernels of $\boldsymbol\Gamma_\pi(U'^{\tilde u})$ and $\boldsymbol\Gamma_\pi(U^\natural)$ on $\Lam_{M'}(z)$ differ by at most $C\sum_{w}e^{-\gamma(|y-w|+|w-y'|)}$, summed over $w$ within $10M$ of the complement of $\Lam_{20M'}(z)$. These $w$ are at distance at least $9M'$ from $\Lam_{M'}(z)$, which gives the error $Ce^{-\gamma M'}$ after a Schur bound.

By Part~II at the background $U^\natural$, with the averaging pair $(Q'_k(U^\natural),\bQ_k(U^\natural))$ (Remark~\ref{rem:general-averaging}), $\boldsymbol\Gamma_\pi(U^\natural)-\boldsymbol\Gamma(1)$ has decay $(C_{M'}\eps,\gamma)$. This follows from Theorem~\ref{thm:cov-min}(a) and Lemma~\ref{lem:perturbed-inverse}.

Finally, by Lemma~\ref{lem:avg-local-gauge}, whose hypotheses are the third condition defining $\eps_{\rm cmp}(M')$, $W^{\tilde u}=e^{\ii B_W}$ near $\Lam_{M'}(z)$ with $|B_W|\le C_{M'}\eps$. So the transports in $l'(W^{\tilde u})$ are within $C_{M'}\eps$ of $1$ there. Together with Lemma~\ref{lem:one-step-expansion}(b), this gives $|l'(W^{\tilde u})B^u-qB^u|\le C_{M'}\eps\|B\|$.
\end{proof}

\begin{theorem}[Coercivity on the constraint set]\label{thm:step-coercive}
There are $c_L>0$, depending only on $d$ and $L$, and $M'_L\ge M$ and $\eps_L>0$, depending only on $d$, $L$, $G$, $a$, $\alpha$ and $M$, such that, if $\eps\le\eps_L$ and $n\ge n_0$ with $n_0\ge80M'_L$,
\[
\ip B{(\boldsymbol\Gamma_\pi-a)B}\ge c_L\|B\|^2\qquad\text{for all real }B\in\mathcal Z.
\]
\end{theorem}

\begin{proof}
Let $\{\chi_j\}$ be the partition of Lemma~\ref{lem:partition} at scale $M'$, restricted to the unit lattice, acting on bond fields through $b_-$. The operator $K=\boldsymbol\Gamma_\pi-a$ is bounded with decay $(C,\gamma)$ (Theorem~\ref{thm:reg-min}). By Lemma~\ref{lem:IMS} and the argument of Lemma~\ref{lem:double-comm} for kernels, $\sum_j(\chi_j(x)-\chi_j(x'))^2|K(x,x')|\le CM'^{-2}|x-x'|^2|K(x,x')|$, so
\[
\ip BKB\ge\sum_j\ip{\chi_jB}{K\chi_jB}-CM'^{-2}\|B\|^2 .
\]
Each $\chi_jB$ vanishes on $\mathcal T$ and is supported on bonds starting in a cube of side $2M'$. By Lemma~\ref{lem:local-comparison}, Proposition~\ref{prop:eff-action-gauge} and Lemma~\ref{lem:poincare-step},
\[
\begin{aligned}
\ip{\chi_jB}{K\chi_jB}&\ge c_D\|\dl(\chi_jB)^{u_j}\|^2-\delta_1\|\chi_jB\|^2\\
&\ge c_DC_{\rm P}^{-1}\|\chi_jB\|^2-c_D\|q(\chi_jB)^{u_j}\|^2-\delta_1\|\chi_jB\|^2,
\end{aligned}
\]
with $\delta_1=C_{M'}\eps+Ce^{-\gamma M'}$. Moreover $\|q(\chi_jB)^{u_j}\|^2\le2\|l'\chi_jB\|^2+2C_{M'}^2\eps^2\|\chi_jB\|^2$. The operator $l'$ has range $2L$, and $|\chi_j(x)-\chi_j(x')|\le CL/M'$ within this range. With the bounded overlap of the $\chi_j$, this gives $\sum_j\|l'\chi_jB\|^2\le2\|l'B\|^2+C(L/M')^2\|B\|^2$. For $B\in\mathcal Z$, $l'B=0$. Summing over $j$ with $\sum_j\|\chi_jB\|^2=\|B\|^2$,
\[
\ip BKB\ge\Big(c_DC_{\rm P}^{-1}-CM'^{-2}-Cc_DL^2M'^{-2}-Ce^{-\gamma M'}-C_{M'}\eps\Big)\|B\|^2 .
\]
Choose $M'=M'_L$ so large that the terms independent of $\eps$ are at most $\frac14c_DC_{\rm P}^{-1}$. Then choose $\eps_L\le\eps_{\rm cmp}(M'_L)$, below the threshold of Corollary~\ref{cor:parts-apply} for Theorem~\ref{thm:reg-min} (which gives the decay of $K$ used in the IMS step), so small that $C_{M'_L}\eps_L\le\frac14c_DC_{\rm P}^{-1}$, and put $c_L=\frac12c_DC_{\rm P}^{-1}$. Lemma~\ref{lem:local-comparison} applies, since $n\ge80M'_L$ and $\eps\le\eps_{\rm cmp}(M'_L)$.
\end{proof}

\subsection{The fluctuation covariance}

\begin{corollary}[Fluctuation covariance]\label{cor:covariance}
Let $\eps\le\eps_L$ and $CL^{3d+1}\eps\le\frac14c_L$. Then $\Delta_k$, as an operator on real fields on $\mathfrak F$, is symmetric with
\[
\tfrac14c_L\|Z\|^2\le\Delta_k(Z,Z)\le C_L\|Z\|^2 .
\]
It has decay $(C_L,\gamma_1)$, and $C_k=\Delta_k^{-1}$ has decay $(C'_L,\gamma_C)$. For complex base configurations $\hat W=e^{\ii\hat\beta}W$ with $|\hat\beta|_\infty\le c'_L$, where $c'_L\le r_1/4$ is a small constant, define
\[
\Delta_k(\hat W)(Z,Z)=\Delta^0(\hat W)\big(\hat Z^1_{\hat W}Z,\hat Z^1_{\hat W}Z\big)+2\ip{\omega_0(\hat W)}{\hat Z^2_{\hat W}(Z)},
\]
with the objects of Proposition~\ref{prop:base-dependence} and the elimination of Lemma~\ref{lem:eliminate} at the base $\hat W$. Then:
\begin{itemize}
\item $\Delta_k(\hat W)$ and $C_k(\hat W)$ are analytic in $\hat\beta$;
\item $\Rea\Delta_k(\hat W)(Z,\bar Z)\ge\frac18c_L\|Z\|^2$;
\item the decay bounds persist, and $|\partial_{\hat\beta(c)}\Delta_k(\hat W)(b,b')|+|\partial_{\hat\beta(c)}C_k(\hat W)(b,b')|\le C_Le^{-\gamma_C\tau(c,b,b')}$.
\end{itemize}
\end{corollary}

\begin{proof}
By Proposition~\ref{prop:hessian}(b) and Theorem~\ref{thm:step-coercive}, $\Delta^0(\hat Z^1Z,\hat Z^1Z)\ge(c_L-C\eps)\|\hat Z^1Z\|^2\ge\frac12c_L\|Z\|^2$, since $\hat Z^1Z\in\mathcal Z$ and $\|\hat Z^1Z\|\ge\|Z\|$. By Lemma~\ref{lem:eliminate}(d) and $|\omega_0|\le C\eps$, $|2\ip{\omega_0}{\hat Z^2(Z)}|\le CL^{3d+1}\eps\|Z\|^2$. The upper bound and the decay follow from Proposition~\ref{prop:hessian}(c) and the locality of $\hat Z^1$, $\hat Z^2$. Lemma~\ref{lem:CT-bounded} gives the decay of $C_k$.

For complex bases, the elimination of Lemma~\ref{lem:eliminate} is analytic in $\hat\beta$. The function $\ell_c$ is an analytic function of the finitely many bond variables. The proof of Lemma~\ref{lem:one-step-expansion} used unitarity of $R(W(\cdot))$ only to bound transports along contours of at most $3dL$ bonds, and these remain bounded by $1+CL|\hat\beta|_\infty$ for complex bases, which is harmless for $c'_L$ small. Together with Proposition~\ref{prop:base-dependence}, this shows that $\Delta_k(\hat W)-\Delta_k(W)$ has decay $(C_L|\hat\beta|_\infty,\gamma_1)$ and derivatives with tree decay. The lower bound follows for $c'_L$ small. Lemma~\ref{lem:CT-bounded}, in the complex form of Section~\ref{sec:perturbation}, gives the decay of $C_k(\hat W)$. Finally $\partial C_k=-C_k(\partial\Delta_k)C_k$, and $\tau(c,b_1,b_2)+|b-b_1|+|b_2-b'|\ge\tau(c,b,b')$ gives the tree bound.
\end{proof}

\begin{remark}[The constants do not depend on $M$]\label{rem:step-M}
The constants $c_L$, $C_L$, $C'_L$, $\gamma_1$, $\gamma_C$ and $c'_L$ of Theorem~\ref{thm:step-coercive} and Corollary~\ref{cor:covariance}, and the constant $C_L$ of Theorem~\ref{thm:step}, do not depend on $M\ge M_0$:
\begin{itemize}
\item $c_L=\frac12c_DC_{\rm P}^{-1}$, with $c_D$ of Proposition~\ref{prop:eff-action-gauge} and $C_{\rm P}$ of Lemma~\ref{lem:poincare-step}; they depend only on $d$ and $L$.
\item The upper bound and the decay $(C_L,\gamma_1)$ of $\Delta_k$ come from Proposition~\ref{prop:hessian}(c), that is, Proposition~\ref{prop:coarse-action}(c) at $B=0$, and from the locality of $\hat Z^1$ and $\hat Z^2$. The constants of Proposition~\ref{prop:coarse-action} come from Theorem~\ref{thm:variational}, whose constants may be taken independent of $M\ge M_0$ (the paragraph after the theorem), from Propositions~\ref{prop:action} and~\ref{prop:Q-expansion} (the first assumes only $(1,\eps)$-regularity; the proofs of Sections~\ref{sec:averaging} and~\ref{sec:adapted} do not use $M$, see Section~\ref{sec:AB-proof}, Setting; Lemma~\ref{lem:eliminate} depends only on $d$ and $G$), and from $r_1=c_6/4$. The comparison of Proposition~\ref{prop:hessian}(b) uses Lemma~\ref{lem:X-coercive}, whose constant comes from Theorem~\ref{thm:reg-coercive} and does not depend on $M\ge M_0$.
\item $C'_L$ and $\gamma_C$ come from these by Lemma~\ref{lem:CT-bounded}. The radius $c'_L\le r_1/4$ and the bounds at complex bases, through Proposition~\ref{prop:base-dependence}, are of the same kind.
\item The constant $C_L$ of Theorem~\ref{thm:step}(d) comes from the proof of Lemma~\ref{lem:step-quadratic}, on the disc $|Z|_\infty\le2cL^{-2d-2}$, where $c$ depends only on $d$, $G$, $L$ and $c_6$, none of which involves $M$.
\end{itemize}
The parameter $M$ enters this part only through thresholds and radii: through $M'_L\ge M$; through the constant $C_{M'}$ and the factor $e^{CM}$ of Theorem~\ref{thm:reg-locality}(b) in Lemma~\ref{lem:local-comparison}, which are absorbed by the choice of $M'_L$ and of $\eps_L\le\eps_{\rm cmp}(M'_L)$; and through $a'_1$, $a_{\rm step}$ and $r_2$, which involve the constants $a_1$, $B_3$ of Theorem~\ref{thm:balaban} and $a_\delta$ of Proposition~\ref{prop:sharp-curvature}.
\end{remark}

The fluctuation field lives on the unit lattice $\T$, where all operators are bounded. So the decay of the covariance follows from Combes--Thomas directly. In particular we do not need Ba\l aban's representation of the axial-gauge covariance through Landau-gauge operators on the fine lattice \cite[(3.183)--(3.185)]{BalBackground}. The fine-lattice Green functions of Parts~I--III enter only through $\Delta^0$, i.e.\ through the coarse effective action.

\subsection{The determinant}

The Gaussian integral over the free variables produces the factor $\det(\Delta_k)^{-1/2}$, which depends on the background. We write its logarithm as a sum of local densities.

\begin{proposition}[Determinant densities]\label{prop:det-densities}
For real bases let
\[
\log\Delta_k=\int_0^\infty\Big(\frac1{1+t}-(\Delta_k+t)^{-1}\Big)dt,\qquad\ell_k(b)=\operatorname{tr}_{\mathfrak g}\big[(\log\Delta_k)(b,b)\big],\quad b\in\mathfrak F .
\]
Then $\log\det\Delta_k=\sum_{b\in\mathfrak F}\ell_k(b)$, and:
\begin{enumerate}[label=(\alph*)]
\item $\log\Delta_k$ has decay $(C_L,\gamma_C/2)$; in particular $|\ell_k(b)|\le C_L$;
\item $\ell_k(b)$ is invariant under real gauge transformations of $(W,U')$;
\item for complex bases $\hat W=e^{\ii\hat\beta}W$ as in Corollary~\ref{cor:covariance}, the same integral defines $\ell_k(b)$ analytically, and $|\partial_{\hat\beta(c)}\ell_k(b)|\le C_Le^{-\gamma_C|b-c|/2}$;
\item at the trivial base $W=1$, $\ell_k(b)$ depends only on the position of $b$ relative to the blocks $B'(\cdot)$.
\end{enumerate}
\end{proposition}

\begin{proof}
(a) For $t\ge0$, $\Delta_k+t$ has decay $(C_L+t,\gamma_1)$ and is bounded below by $\frac14c_L+t$. By Lemma~\ref{lem:CT-bounded} its inverse has decay $(C(c_L/4+t)^{-1},\gamma_C)$, with $\gamma_C$ independent of $t$, since the ratio of the lower bound to the decay constant does not degenerate as $t\to\infty$. The integrand equals $(\Delta_k-1)(1+t)^{-1}(\Delta_k+t)^{-1}$, which has decay $(C_L(1+t)^{-1}(c_L/4+t)^{-1},\gamma_C/2)$ and is integrable. For a positive symmetric matrix the integral is the logarithm, and $\operatorname{Tr}\log=\log\det$.

(b) The fluctuation coordinates are covariant: a real gauge transformation $u$ acts on fields on $\mathfrak F$ bondwise by $R(u(b_-))$, which is orthogonal. By Proposition~\ref{prop:hessian}(d) and the covariance of the elimination, $\Delta_k\mapsto R(u)\Delta_kR(u)^{-1}$. Hence $\log\Delta_k$ is conjugated in the same way, and its diagonal traces are invariant.

(c) For complex bases, $\Rea(\Delta_k+t)\ge\frac18c_L+t$, so the integral converges. It is analytic and agrees with the logarithm for real bases. The derivative is $\int_0^\infty(\Delta_k+t)^{-1}(\partial\Delta_k)(\Delta_k+t)^{-1}dt$. By the tree bound of Corollary~\ref{cor:covariance} for $\partial\Delta_k$, its diagonal kernel at $b$ is bounded by $C_L\sum_{b_1,b_2}e^{-\gamma_C(|b-b_1|+|b_2-b|)}e^{-\gamma_1\tau(c,b_1,b_2)}(c_L/8+t)^{-2}$, which integrates to $C_Le^{-\gamma_C|b-c|/2}$.

(d) At $W=1$ the operator $\Delta_k$ commutes with the translations by $L\Z^d$.
\end{proof}

Proposition~\ref{prop:det-densities} is the starting point of the analysis of Ba\l aban's normalization factor \cite[(61)--(63)]{BalUV3}. What remains is to show that $\ell_k(b;W)-\ell_k(b;1)$ is a gauge-invariant function of the background near $b$, up to exponentially small terms, and that it is small. This requires the comparison of $\Delta_k$ at bases that differ by a large amount far from $b$. It is part of the localization of the background dependence. This is carried out in Part~VIII (Section~\ref{sec:det-local}) with Dirichlet approximants, and in Part~IX with weakened random walks.

\subsection{One step as a Gaussian integral}

\begin{theorem}[One small-field step]\label{thm:step}
There are $a_{\rm step}>0$ and $r_2>0$, depending only on $d$, $L$, $G$, $a$, $\alpha$ and $M$, and $C_L$, depending only on $d$, $L$, $G$, $a$ and $\alpha$, with the following property. Let $k\ge1$. Let $V'$ be real on $\T'$ with curvature at most $\eps'_1\le a_{\rm step}$ in the units of step $k+1$. Let $W=V^0$, $U'$ and the fluctuation coordinates be as above, and let $w\le r_2$. Let $\rho$ be a gauge-invariant density of the form \eqref{eq:rho-form} on $\mathfrak U_k(\tilde h)$, with $a_1\ge\tilde h\ge2(C_*+1)L^{-2}\eps'_1+Cw$, and let $g=g_k$. Then
\[
\begin{aligned}
T(\chi^{\rm w}\rho)(V')={}&\mathcal N_k\exp\Big(-\frac1{g^2}S(U')-\frac12\sum_{b\in\mathfrak F}\ell_k(b)+\sum_c\log\mathfrak j_c\Big)\\
&\times\int d\mu_{C_k}(Y)\ \chi^{\rm w}\big(\hat Z(gY)\big)\exp\Big(-\frac1{g^2}\mathcal V(gY)+\mathfrak J(gY)+\mathfrak E\big(e^{\ii\hat Z(gY)}W\big)\Big),
\end{aligned}
\]
where the integrand vanishes unless $g|Y(b)|\le w$ for all $b\in\mathfrak F$, as in Proposition~\ref{prop:exact-step}, and:
\begin{enumerate}[label=(\alph*)]
\item $\mathcal N_k=(2\pi g^2)^{|\mathfrak F|\dim\mathfrak g/2}\,\sigma(0)^{|\mathfrak F|}\,L^{(d-1)\dim\mathfrak g\,|\mathcal B(\T')|}$ does not depend on the background; $|\mathcal B(\T')|$ is the number of bonds of $\T'$;
\item $\mathfrak j_c=L^{-(d-1)\dim\mathfrak g}\mathsf J_c(0)=1+O(L^d\mathsf a)$ is a local function of the base near $c$;
\item $d\mu_{C_k}$ is the Gaussian measure with the covariance $C_k$ of Corollary~\ref{cor:covariance};
\item $\mathcal V$ is analytic for complex $|Z|_\infty\le r_2$, of third order, with $|g^{-2}\mathcal V(gY)|\le C_Lg|Y|_\infty\|Y\|^2$;
\item $\mathfrak J(Z)=\sum_{b\in\mathfrak F}\log(\sigma(Z(b))/\sigma(0))+\sum_c\log(\mathsf J_c(Z)/\mathsf J_c(0))$ is a sum of local analytic terms, bounded by $C|Z(b)|^2$ and by $CL^{2d+1}|Z|_\infty$ respectively, with $|Z|_\infty$ over the free bonds (equivalently by $CL^{d+1}w$ on the support of $\chi^{\rm w}$, where all $|\hat Z(b)|\le w$);
\item $\mathfrak E(e^{\ii\hat Z(Z)}W)=\mathfrak E\big(e^{\ii\eta A_*(\hat Z(Z))}U'\big)$, where $\mathfrak E(U)$ denotes the gauge-invariant function $\mathfrak E_k$ expressed through the background field.
\end{enumerate}
The term $g_k^{-2}S(U')$ in the units of step $k$ is $g_{k+1}^{-2}S(U')$ in the units of step $k+1$.
\end{theorem}

\begin{proof}
Choose $a_{\rm step}$ so small that $4(C_*+1)L^{-2}a_{\rm step}\le a'_1$, and that Lemma~\ref{lem:composition}, Lemma~\ref{lem:eliminate} (with $L^d\mathsf a\le c$), Theorem~\ref{thm:step-coercive} and Corollary~\ref{cor:covariance} apply with $\eps=B_3\eps'_1$. By Proposition~\ref{prop:exact-step}, $T(\chi^{\rm w}\rho)(V')$ is an integral over $Z$ on $\mathfrak F$. On the support of $\chi^{\rm w}$, which bounds all bonds including $b_0(c)$, the configuration $e^{\ii\hat Z(Z)}W$ has curvature at most $2(C_*+1)L^{-2}\eps'_1+Cw\le\tilde h$. So \eqref{eq:rho-form} applies, and $\rho(e^{\ii\hat Z}W)=\exp(-g^{-2}F(Z)+\mathfrak E)$ with $F$ as in Lemma~\ref{lem:step-quadratic}. Since $\mathfrak j_c=L^{-(d-1)\dim\mathfrak g}\mathsf J_c(0)$ by~(b), write $\mathsf J_c(Z)=L^{(d-1)\dim\mathfrak g}\mathfrak j_c\cdot\mathsf J_c(Z)/\mathsf J_c(0)$ and $\sigma(Z(b))=\sigma(0)\,\sigma(Z(b))/\sigma(0)$. The Gaussian integral
\[
\int\prod_{b\in\mathfrak F}dZ(b)\,e^{-\frac1{2g^2}\Delta_k(Z,Z)}f(Z)=(2\pi g^2)^{|\mathfrak F|\dim\mathfrak g/2}(\det\Delta_k)^{-1/2}\int d\mu_{C_k}(Y)f(gY)
\]
and Proposition~\ref{prop:det-densities} give the formula. (d) follows from Lemma~\ref{lem:step-quadratic}. (e) follows from $\sigma(Z)/\sigma(0)=1+O(|Z|^2)$ and Proposition~\ref{prop:exact-step}. (f) is gauge invariance: by Theorem~\ref{thm:variational}(c), $U_k(e^{\ii\hat Z}W)$ and $e^{\ii\eta A_*(\hat Z)}U'$ differ by a real gauge transformation of $\Tf$, which need not lie in $\mathfrak G_0$.
\end{proof}

\begin{remark}[The first step]\label{rem:step-k0}
Theorem~\ref{thm:step} is stated for $k\ge1$, because Sections~\ref{sec:coarse-action} and~\ref{sec:covariance} use the $k$-step operators of Parts~II--V. For $k=0$ the same representation holds, with its own constants $a^{(0)}_{\rm step}$, $r^{(0)}_2$ and $C_L$, depending on the same quantities. At $k=0$, $\eta=1$, $\T=\Tf$ and $U_0(V)=V$ (Theorem~\ref{thm:balaban}), so $W=U'^{(0)}=U'$, $A_*(B)=B$ and $\mathcal S(B)=S(e^{\ii B}U')$. The proof of Theorem~\ref{thm:step} applies with the following inputs.
\begin{itemize}
\item Lemma~\ref{lem:composition}, Lemma~\ref{lem:eliminate} and Proposition~\ref{prop:exact-step} hold at $k=0$, with $V^0=U'$: the $0$-fold average is the identity, and the other statements concern the one-step averaging $\mathfrak a'$ at the base $W$ and Theorem~\ref{thm:balaban} for $k=0$ and $k=1$. For Lemma~\ref{lem:composition}(c) at $k=0$: a real $V\in\mathfrak U_0(a'_1)$ with $\mathfrak a'(V)=V'$ satisfies $U_0(V)=V\in\mathfrak U(B_3a'_1)$, by Theorem~\ref{thm:balaban}(b) at $k=0$. In the units of step $1$ its curvature and current are at most $L^3B_3a'_1\le\theta_0$, so $V\in\mathfrak B_1(V')\cap\mathfrak U(\theta_0)$, and Theorem~\ref{thm:balaban}(a) at step $1$ gives $S(V)\ge S(U')$, with equality only on the orbit of $U'$, as in the proof of Lemma~\ref{lem:composition}(c).
\item Lemma~\ref{lem:step-quadratic} holds with this $\mathcal S$, with $\Delta^0$ the Hessian of $B\mapsto S(e^{\ii B}U')$ at $0$ and $\omega_0=J_{U'}$. The linear term vanishes by Lemma~\ref{lem:composition}(c) at $k=0$. The cubic bound $|\mathcal V(Z)|\le C_L|Z|_\infty\|Z\|^2$ follows from Cauchy's estimate for the plaquette terms of $S$, each of which depends on $Z$, through $\hat Z$, on at most $C_L$ bonds within distance $2L+2$ of the plaquette.
\item Theorem~\ref{thm:step-coercive} is not used. The lower bound $\Delta_0\ge\frac12C_{\rm P}^{-1}$ of Proposition~\ref{prop:first-step}(b) replaces $\frac14c_L$ in Corollary~\ref{cor:covariance} and in the proof of Proposition~\ref{prop:det-densities}. Its proof treats the local operator $\Delta^0$ directly and does not use Theorem~\ref{thm:step}. The decay of $\Delta_0$ and of $C_0=\Delta_0^{-1}$ follows from the locality of $\Delta_0$ and Lemma~\ref{lem:CT-bounded}.
\item The constant $a^{(0)}_{\rm step}$ is chosen so that Lemma~\ref{lem:composition}, Lemma~\ref{lem:eliminate} (with $L^d\mathsf a\le c$) and $4(C_*+1)L^{-2}\eps'_1\le a'_1$ hold, together with the smallness conditions on $\eps=B_3\eps'_1$ in the proof of Proposition~\ref{prop:first-step}(b), among them $\eps\le\eps_{\rm cmp}(M')$. The radius $r^{(0)}_2\le r_L$ is chosen as $r_2$ in Lemma~\ref{lem:step-quadratic}, with $r_1$ replaced by $1$.
\end{itemize}
The formula of Theorem~\ref{thm:step} then holds at $k=0$, with $g=g_0$, $\mathfrak E=\mathfrak E_0$, and $\mathcal N_0$, $\mathfrak j_c$ and $\mathfrak J$ as in (a), (b) and (e).
\end{remark}

\begin{remark}[What Parts~VIII--X have to do]
Theorem~\ref{thm:step} is exact on the small-field set. The analysis of its Gaussian integral is the content of Parts~VIII--X:
\begin{itemize}
\item a cluster expansion, in Part~X, of the Gaussian integral of Theorem~\ref{thm:step} (with the cutoff $\chi^{\rm w}(\hat Z(gY))$ and the exponent $-g^{-2}\mathcal V(gY)+\mathfrak J(gY)+\mathfrak E(e^{\ii\hat Z(gY)}W)$), which requires a localized (random-walk) representation of $C_k$ and of the coefficients of $\mathcal V$;
\item the localization of the dependence on the background of the densities $\ell_k(b)$, of $\mathfrak j_c$ and of the new activities;
\item the normalization at $U'=1$, which removes the vacuum energy, and the proof that the remainder is of second order in the curvature of the background.
\end{itemize}
The inputs proved here are the coercivity and decay of $C_k$ (Corollary~\ref{cor:covariance}), the tree decay of the derivatives of the coarse action (Proposition~\ref{prop:coarse-action}) and of its dependence on the base (Proposition~\ref{prop:base-dependence}), and the analyticity of all ingredients in a complex neighbourhood.
\end{remark}
